\chapter{线性注意力的挑战与改进}

\section{表达能力的差距}

线性注意力面临的最主要挑战是表达能力不如标准 softmax 注意力。

\subsection{问题分析}

标准 softmax 注意力能够学习尖锐的注意力分布，即对少数关键位置赋予高权重，对其他位置赋予低权重。这是因为 softmax 的指数放大效应：

\begin{equation}
    \text{softmax}(s_i) = \frac{\exp(s_i)}{\sum_j \exp(s_j)}
\end{equation}

当某个 $s_i$ 显著大于其他值时，softmax 会赋予其接近 1 的权重。这种"聚焦"能力对许多任务至关重要。

相比之下，线性注意力使用的核函数（如 elu+1、随机特征）无法产生同样尖锐的分布。具体而言：

\begin{itemize}
    \item elu+1 映射产生的核值范围有限，无法像指数函数那样放大差异
    \item 随机特征近似引入了近似误差，尤其在极端值情况下
\end{itemize}

\subsection{实验观察}

研究者发现线性注意力在以下场景下与标准注意力差距较大：

\begin{itemize}
    \item 需要精确检索的任务（如 copying task）
    \item 需要关注长距离依赖的任务
    \item 需要尖锐注意力分布的任务
\end{itemize}

\section{数值稳定性问题}

线性注意力在训练中可能遇到数值稳定性问题。

\subsection{分母接近零}

在归一化公式中：

\begin{equation}
    \mathbf{o}_i = \frac{\phi(\mathbf{q}_i)^\top \mathbf{S}}{\phi(\mathbf{q}_i)^\top \mathbf{z}}
\end{equation}

如果分母 $\phi(\mathbf{q}_i)^\top \mathbf{z}$ 接近零，会导致数值不稳定。实践中通常添加一个小常数 $\epsilon$：

\begin{equation}
    \mathbf{o}_i = \frac{\phi(\mathbf{q}_i)^\top \mathbf{S}}{\phi(\mathbf{q}_i)^\top \mathbf{z} + \epsilon}
\end{equation}

但这只是部分缓解，不能完全解决问题。

\subsection{状态矩阵的累积}

在递推形式中，状态矩阵 $\mathbf{S}_t = \sum_{j=1}^t \phi(\mathbf{k}_j)\mathbf{v}_j^\top$ 会随着序列长度增长而累积，可能导致数值溢出或精度损失。

\section{改进方法}

\subsection{混合注意力}

一种实用的策略是混合使用标准注意力和线性注意力：

\begin{equation}
    \mathbf{O} = \lambda \cdot \text{SoftmaxAttn}(\mathbf{Q}, \mathbf{K}, \mathbf{V}) + (1-\lambda) \cdot \text{LinearAttn}(\mathbf{Q}, \mathbf{K}, \mathbf{V})
\end{equation}

或者将序列分为两部分：局部窗口内使用标准注意力，全局使用线性注意力。

\subsection{改进的特征映射}

\subsubsection{CosFormer}

CosFormer 使用余弦函数来重构特征映射：

\begin{equation}
    \phi(\mathbf{x}) = \text{relu}(\mathbf{x}) \odot \cos\left(\frac{\pi \cdot \text{pos}}{M}\right)
\end{equation}

其中位置编码通过余函数引入，使得近位置的核值更大，远位置自然衰减。

\subsubsection{NormFormer}

NormFormer 通过输出归一化来稳定训练：

\begin{equation}
    \mathbf{O} = \text{LayerNorm}\left(\text{LinearAttn}(\mathbf{Q}, \mathbf{K}, \mathbf{V})\right)
\end{equation}

\subsection{训练策略改进}

\subsubsection{课程学习}

一种有效的训练策略是：

\begin{enumerate}
    \item 先用标准 softmax 注意力预训练模型
    \item 然后替换为线性注意力，用较小的学习率微调
    \item 逐渐增加序列长度
\end{enumerate}

这种方法可以让模型先在标准注意力下学习良好的表示，再适应线性注意力的近似。

\subsubsection{蒸馏}

使用标准注意力模型作为教师模型，线性注意力模型作为学生模型进行知识蒸馏：

\begin{equation}
    \mathcal{L} = \alpha \mathcal{L}_{\text{task}} + (1-\alpha) \text{KL}(\mathbf{O}_{\text{teacher}} \| \mathbf{O}_{\text{student}})
\end{equation}

\section{理论分析}

\subsection{近似误差界}

设标准注意力输出为 $\mathbf{O}^*$，线性注意力输出为 $\hat{\mathbf{O}}$，则近似误差可以界定为：

\begin{equation}
    \|\mathbf{O}^* - \hat{\mathbf{O}}\|_F \leq C \cdot \max_i \|\phi(\mathbf{q}_i) - \phi^*(\mathbf{q}_i)\| \cdot \|\mathbf{V}\|_F
\end{equation}

其中 $\phi^*$ 是理想的特征映射，$C$ 是常数。

\subsection{表达能力分析}

有研究表明，线性注意力的表达能力受限于特征映射的秩。具体地，线性注意力可以表示的函数空间维度为 $\mathcal{O}(md)$，其中 $m$ 是特征维度，$d$ 是隐藏维度。

相比之下，标准注意力可以表示的函数空间维度为 $\mathcal{O}(Nd)$，当 $N \gg m$ 时，标准注意力具有更强的表达能力。

\section{实践建议}

根据现有的研究和实践经验，以下建议可以帮助更好地使用线性注意力：

\begin{enumerate}
    \item 对于需要长序列处理的场景，线性注意力是可行的选择
    \item 使用混合注意力或局部窗口策略来弥补表达能力的不足
    \item 在训练时添加梯度裁剪和层归一化来提高稳定性
    \item 特征映射的选择应根据任务特性调整：elu+1 简单高效，随机特征更接近 softmax
    \item 推理时利用递推形式实现 $\mathcal{O}(1)$ 每 token 的高效生成
\end{enumerate}
